\section{Acondicionamiento, instrumentaci\'on y digitalizaci\'on}
\label{sec:acondicionamiento-digitalizacion}

Esta secci\'on comienza en los terminales el\'ectricos de entrada del nodo. No vuelve a
discutir qu\'e transductor se emplea, su modelo mec\'anico ni la fuente que genera la
se\~nal. El problema acotado es otro: \textbf{convertir una tensi\'on diferencial peque\~na
en una secuencia digital que aproveche el rango del conversor, conserve la informaci\'on
de fase y pueda adquirirse de forma determinista}. La transmisi\'on de esas muestras se
trata despu\'es, como una etapa independiente.

El desarrollo de esta parte no produjo un circuito definitivo de una sola vez. El
historial de Git del antiguo repositorio \texttt{tesis-legacy}, preservado actualmente en
\texttt{src/firmware/psoc}, permite reconstruir una sucesi\'on de proyectos de PSoC
Creator entre el 17 de abril y el 3 de junio de 2026, seguida por revisiones de
calibraci\'on, DMA y control de captura hasta agosto. Esa evoluci\'on es relevante porque
explica qu\'e problema resolvi\'o cada bloque que permanece en el dise\~no.

\subsection{Qu\'e debe resolver esta etapa}

Si $v_d(t)$ es la tensi\'on diferencial disponible en la entrada, una representaci\'on
compacta del camino de acondicionamiento y conversi\'on es

\begin{equation}
v_{\mathrm{ADC}}(t)=V_{\mathrm{CM}}+
\left[h_{\mathrm{AFE}}(t)*v_d(t)\right]+v_{\mathrm{os}}(t)+n_{\mathrm{AFE}}(t),
\label{eq:afe-general}
\end{equation}

donde $H_{\mathrm{AFE}}(s)$ incluye la ganancia y los filtros anal\'ogicos,
$V_{\mathrm{CM}}$ fija el punto de operaci\'on para alimentaci\'on simple,
$v_{\mathrm{os}}$ agrupa los desplazamientos de continua y $n_{\mathrm{AFE}}$ el ruido
incorporado por la instrumentaci\'on. El ADC observa esta tensi\'on en $t=nT_s$ y entrega

\begin{equation}
x[n]=Q_{B,V_{\mathrm{span}}}\!\left\{v_{\mathrm{ADC}}(nT_s)\right\},
\qquad
\Delta V_{\mathrm{LSB}}=\frac{V_{\mathrm{span}}}{2^B}.
\label{eq:adc-cuantizacion}
\end{equation}

La Ec.~\eqref{eq:afe-general} muestra que ``amplificar'' no es suficiente. La etapa debe
resolver simult\'aneamente:

\begin{itemize}
    \item \textbf{entrada diferencial y rechazo de modo com\'un}, para que interferencias
    presentes de forma semejante en ambos conductores no se confundan con la se\~nal;
    \item \textbf{ganancia programable}, porque una ganancia peque\~na desperdicia
    resoluci\'on y una demasiado alta recorta la forma de onda;
    \item \textbf{conformaci\'on espectral anal\'ogica}, incluida la red de compensaci\'on
    y un pasa-bajos anterior al muestreo;
    \item \textbf{control del punto de operaci\'on}, ya que el \emph{offset} de varias
    etapas en cascada puede consumir el rango del ADC aun sin se\~nal;
    \item \textbf{digitalizaci\'on reproducible}, con tasa, rango, filtrado y formato de
    muestra conocidos;
    \item \textbf{movimiento determinista de datos}, evitando que la atenci\'on de otras
    tareas del procesador cambie el instante de adquisici\'on.
\end{itemize}

Por eso acondicionamiento y digitalizaci\'on se presentan juntos: el rango del ADC decide
cu\'anta ganancia anal\'ogica conviene, el filtro anti-\emph{aliasing} debe actuar antes
del ADC y el multiplexor del propio camino digital permite observar y ajustar las etapas
anal\'ogicas.

\subsection{Reconstrucci\'on mediante Git: todos los proyectos de la cadena GEO}
\label{sec:git-topdesign}

Se buscaron en todas las ramas los archivos binarios
\texttt{TopDesign/TopDesign.cysch} y se sigui\'o su aparici\'on, renombrado y eliminaci\'on.
Las Tablas~\ref{tab:git-circuitos-a} y~\ref{tab:git-circuitos-b} enumeran
\textbf{todos los proyectos de acondicionamiento GEO encontrados}; no se cuentan los
ejemplos USB de Cypress ni los proyectos del martillo, porque pertenecen a otra cadena
instrumental.

\begin{table}[ht]
\centering
\scriptsize
\caption{Primera mitad del inventario hist\'orico de circuitos PSoC Creator.}
\label{tab:git-circuitos-a}
\begin{tabularx}{\textwidth}{@{}>{\raggedright\arraybackslash}p{0.29\textwidth} p{0.12\textwidth} Y@{}}
\toprule
\textbf{Proyecto y primer commit} & \textbf{Fecha} & \textbf{Funci\'on en la evoluci\'on} \\
\midrule
\texttt{SM-24\_V1} --- \texttt{9a91eda} & 17/04 & Primera prueba de acondicionamiento; red pasiva de entrada y bloques b\'asicos de amplificaci\'on. Se abandon\'o al cambiar el reparto entre anal\'ogico y digital. \\
\texttt{Tesis} --- \texttt{74a65c2} & 18/04 & Primera cadena activa temprana, con ramas diferenciales, suma/transimpedancia y filtrado. Reemplaz\'o al proyecto inicial. \\
\texttt{Comparation} --- \texttt{c4b0a8a} & 19/04 & Banco para comparar caminos anal\'ogicos e integrar AMux, ADC $\Delta\Sigma$, filtro digital y salida de laboratorio. \\
\texttt{FiltradoDiferencial} --- \texttt{0051880} & 22/04 & Conserv\'o la se\~nal en forma diferencial a trav\'es de la amplificaci\'on y el filtrado. \\
\texttt{ComparacionDiferencial} --- \texttt{0051880} & 22/04 & Proyecto paralelo para contrastar las variantes diferenciales y su conversi\'on. \\
\texttt{ComparacionDiferencial\_Copy\_01} --- \texttt{c46b7ae} & 30/04 & Rama de comparaci\'on durante el pivote de topolog\'ia; separ\'o etapas PGA, BP y LP para contrastar Sallen--Key/VCVS con MFB. \\
\texttt{DiferencialConRrefs} --- \texttt{6900805} & 10/05 & Introdujo referencias desplazadas por etapa para corregir manualmente el punto de operaci\'on. \\
\bottomrule
\end{tabularx}
\end{table}

\begin{table}[ht]
\centering
\scriptsize
\caption{Segunda mitad del inventario hist\'orico de circuitos PSoC Creator.}
\label{tab:git-circuitos-b}
\begin{tabularx}{\textwidth}{@{}>{\raggedright\arraybackslash}p{0.32\textwidth} p{0.12\textwidth} Y@{}}
\toprule
\textbf{Proyecto y primer commit} & \textbf{Fecha} & \textbf{Funci\'on en la evoluci\'on} \\
\midrule
\texttt{DiferencialConRrefs\_Copy\_01} --- \texttt{6900805} & 10/05 & Copia de trabajo para modificar la arquitectura de referencias sin perder la variante precedente. \\
\texttt{DiferencialConRrefsXD} --- \texttt{1f73657} & 17/05 & Variante experimental del circuito con referencias independientes. \\
\texttt{DiferencialConRrefsYMux} --- \texttt{1f73657} & 17/05 & Agreg\'o multiplexaci\'on para observar nodos intermedios con un mismo ADC. \\
\texttt{DiferencialConRrefsYMux\_Copy\_01} --- \texttt{1f73657} & 17/05 & Ensay\'o la salida single-ended y la sucesi\'on BP--SUM--LP que conducir\'ia a la cadena completa. \\
\texttt{DiferencialToSingleEnded} --- \texttt{610e19a} & 20/05 & Primera implementaci\'on completa: instrumentaci\'on diferencial, conversi\'on single-ended, MFB, sumador, LP, AMux y ADC. \\
\texttt{DiferencialToSingleEnded\_ESP} --- \texttt{70a15d9} & 24/05 & Integraci\'on inicial del circuito completo con el enlace hacia el controlador externo. \\
\texttt{DiferencialToSingleEnded\_\_ESP} --- renombre en \texttt{9e26866} & 25/05 & Continuaci\'on de la integraci\'on; fue absorbido por el proyecto unificado. \\
\texttt{AcondicionamientoAnalogico} --- \texttt{4062297} & 03/06 & Proyecto consolidado y vigente; reuni\'o la cadena anal\'ogica, digitalizaci\'on y l\'ogica de captura. \\
\bottomrule
\end{tabularx}
\end{table}

Los nombres de los proyectos no deben interpretarse como catorce soluciones totalmente
independientes. Git muestra una genealog\'ia: algunas carpetas fueron copias de trabajo,
otras bancos comparativos y otras reemplazos directos. Agrupadas por cambio estructural,
la evoluci\'on puede resumirse as\'i:

\begin{center}
\small
\fbox{pasivo inicial}
$\longrightarrow$
\fbox{cadena activa diferencial}
$\longrightarrow$
\fbox{Sallen--Key vs. MFB}
$\longrightarrow$
\fbox{referencias por etapa + AMux}
$\longrightarrow$
\fbox{diferencial a single-ended}
$\longrightarrow$
\fbox{AFE + ADC + DMA}
\end{center}

Esta lectura evita dos errores: presentar el circuito final como una elecci\'on inmediata
y contar cada autoguardado del mismo \texttt{TopDesign} como una topolog\'ia nueva.

Las tres primeras revisiones conservan, adem\'as del archivo editable de PSoC Creator,
una exportaci\'on completa. Los hitos 02 y 03 ocupan dos hojas cada uno: una para el
camino anal\'ogico y otra para los bloques digitales. Por eso se presentan como cinco
figuras separadas, desde la Fig.~\ref{fig:topdesign-historico-01} hasta la
Fig.~\ref{fig:topdesign-historico-03-digital}, y no como una miniatura conjunta. As\'i
pueden leerse nombres, valores y conexiones sin esconder precisamente la parte de
digitalizaci\'on que diferencia a cada revisi\'on.

\begin{figure}[p]
\centering
\includegraphics[width=\textwidth,height=0.82\textheight,keepaspectratio]{historico_psoc/01_sm24_ina.png}
\caption{Hito 01, \texttt{SM-24\_V1}: primera instrumentaci\'on diferencial
con dos PGA y un operacional de salida. La red de resistencias establece las
referencias y la diferencia entre ramas, pero todav\'ia no aparecen el ADC, el
movimiento de muestras ni un camino de observaci\'on de etapas. La hoja completa
expone el punto de partida real del desarrollo.}
\label{fig:topdesign-historico-01}
\end{figure}

\begin{figure}[p]
\centering
\includegraphics[width=\textwidth,height=0.82\textheight,keepaspectratio]{historico_psoc/02_tesis_analogico.png}
\caption{Hito 02, \texttt{Tesis}, hoja anal\'ogica: el acondicionamiento pasa a dos ramas TIA,
polarizadas alrededor de $V_{DDA}/2$, y termina en un ADC $\Delta\Sigma$ de
17 bits. Los capacitores de realimentaci\'on y las redes RC anteriores al ADC
introducen una conformaci\'on espectral expl\'icita que no exist\'ia en la
revisi\'on inicial.}
\label{fig:topdesign-historico-02-analogico}
\end{figure}

\begin{figure}[p]
\centering
\includegraphics[width=0.72\textwidth,height=0.72\textheight,keepaspectratio]{historico_psoc/02_tesis_digital.png}
\caption{Hito 02, \texttt{Tesis}, hoja digital: filtro digital configurado
como rechazo de banda, transferencia por DMA y una interrupci\'on al guardar
cada muestra. Esta segunda hoja demuestra que la versi\'on no terminaba en el
ADC; ya exist\'ia una ruta determinista de posprocesado y memoria.}
\label{fig:topdesign-historico-02-digital}
\end{figure}

\begin{figure}[p]
\centering
\includegraphics[width=\textwidth,height=0.82\textheight,keepaspectratio]{historico_psoc/03_comparation_analogico.png}
\caption{Hito 03, \texttt{Comparation}, hoja anal\'ogica: banco de comparaci\'on de caminos. A la
izquierda se incorporan un ADC SAR de 12 bits y DMA; a la derecha se conserva el
camino anal\'ogico diferencial, se agregan multiplexores, dos operacionales de
filtrado y un ADC $\Delta\Sigma$ de 18 bits con DMA y filtro digital. La hoja
documenta que la elecci\'on del conversor y el reparto anal\'ogico--digital se
evaluaron como un sistema, no como decisiones independientes.}
\label{fig:topdesign-historico-03-analogico}
\end{figure}

\begin{figure}[p]
\centering
\includegraphics[width=0.82\textwidth,height=0.76\textheight,keepaspectratio]{historico_psoc/03_comparation_digital.png}
\caption{Hito 03, \texttt{Comparation}, hoja digital: filtro y DMA, temporizador
de 16 bits alimentado por el reloj de bus, USBFS y UART a \SI{57600}{bit\per\second}.
Esta hoja completa la arquitectura visible en la figura anterior y registra las
alternativas de temporizaci\'on y comunicaci\'on disponibles en el banco.}
\label{fig:topdesign-historico-03-digital}
\end{figure}

La comparaci\'on binaria de los veinte archivos recuperados tambi\'en encontr\'o un caso
concreto de duplicaci\'on: los \texttt{TopDesign} de
\texttt{FiltradoDiferencial} y \texttt{ComparacionDiferencial} son id\'enticos byte a
byte. Ambos nombres se conservan en el inventario porque corresponden a proyectos del
historial, pero no deben presentarse como dos circuitos distintos. Los otros dieciocho
archivos principales poseen contenido diferente y requieren una captura individual.

\subsection{Decisiones de topolog\'ia que sobreviven en el dise\~no}

\subsubsection{Del filtro pasivo al reparto anal\'ogico--digital}

El primer proyecto intent\'o resolver la selecci\'on de banda mediante una red RC pasiva.
Se descart\'o por sensibilidad a tolerancias, deriva de capacitores y fase no lineal.
En un arreglo multicanal importa no s\'olo que un canal tenga la magnitud esperada, sino
que los canales tengan respuestas de fase reproducibles. El reemplazo estableci\'o una
regla de dise\~no que contin\'ua vigente:

\begin{quote}
Lo que debe ocurrir antes de cuantizar ---ganancia, punto de operaci\'on, compensaci\'on y
anti-\emph{aliasing}--- permanece en anal\'ogico; la selecci\'on reconfigurable de banda y
el movimiento de muestras se realizan con recursos digitales.
\end{quote}

\subsubsection{De Sallen--Key a MFB}

La compensaci\'on adoptada puede escribirse como un camino directo menos un pasa-banda
de segundo orden~\cite{Ma2023}:

\begin{equation}
H_{\mathrm{comp}}(s)=1-
\frac{2(\zeta_1-\zeta_0)\omega_n s}
{s^2+2\zeta_1\omega_n s+\omega_n^2}.
\label{eq:comp-afe}
\end{equation}

Por tanto, el filtro BP y el sumador no son dos bloques arbitrarios: juntos realizan la
resta de la Ec.~\eqref{eq:comp-afe}. Se analizaron Sallen--Key/VCVS y MFB. Las dos
reproduc\'ian de forma semejante la respuesta nominal; la MFB se seleccion\'o por menor
sensibilidad del factor de calidad y de la frecuencia central a las tolerancias, mejor
atenuaci\'on a alta frecuencia y porque sus desventajas habituales ---inversi\'on y baja
impedancia de entrada--- no son decisivas despu\'es de una etapa de baja impedancia.

\begin{figure}[ht]
\centering
\includegraphics[width=0.9\textwidth]{comparacion_topologias.png}
\caption{Comparaci\'on reconstruida a partir del an\'alisis utilizado el 30 de abril.
La magnitud nominal no decide el pivote: el criterio principal es la robustez frente a
tolerancias al replicar canales.}
\label{fig:comparacion-topologias-nueva}
\end{figure}

\begin{figure}[p]
\centering
\includegraphics[width=0.92\textwidth,height=0.76\textheight,keepaspectratio]{mfb_lpf_circuit.png}
\caption{Esquema de la topolog\'ia MFB utilizada durante el an\'alisis de
alternativas. La referencia de alimentaci\'on simple desplaza el punto de trabajo
y las dos redes de realimentaci\'on fijan polos, ganancia y sensibilidad a
tolerancias. La figura permite leer el circuito que est\'a detr\'as de la funci\'on
de transferencia, no s\'olo su diagrama de Bode.}
\label{fig:mfb-circuito}
\end{figure}

\begin{figure}[p]
\centering
\includegraphics[width=\textwidth,height=0.78\textheight,keepaspectratio]{analog_path_check_magcoh.png}
\caption{Comprobaci\'on exploratoria del camino anal\'ogico: magnitud medida,
envolvente de ajuste del circuito, respuesta nominal y coherencia. Se conserva
como evidencia del proceso de diagn\'ostico; la p\'erdida de coherencia a alta
frecuencia impide interpretarla como validaci\'on definitiva de toda la cadena.}
\label{fig:analog-path-check}
\end{figure}

\subsubsection{Conversi\'on diferencial--single-ended y referencias independientes}

La entrada se amplifica con una configuraci\'on de instrumentaci\'on formada por dos PGA
internos del PSoC, coherente con la topolog\'ia documentada por Cypress para sus bloques
de capacitores conmutados~\cite{CypressAN60319}. La cadena se convierte luego a una
se\~nal single-ended referida a un nivel interno. Esto reduce la cantidad de recursos
anal\'ogicos que exigir\'ia mantener dos ramas completas hasta el ADC.

El costo de operar con alimentaci\'on simple y varias etapas de ganancia es la acumulaci\'on
de \emph{offset}. Por eso las versiones \texttt{DiferencialConRrefs*} introdujeron una
referencia independiente para PGA, BP, SUM y LP, y las versiones \texttt{*YMux*} agregaron
un multiplexor que permite medir cada salida intermedia con el mismo ADC. Esa combinaci\'on
es la base de la calibraci\'on asistida digitalmente.

\subsection{Arquitectura anal\'ogica consolidada}
\label{sec:afe-consolidada}

Tomando como entrada $v_d$, la funci\'on de transferencia global puede organizarse como

\begin{equation}
H_{\mathrm{AFE}}(s)=G_{\mathrm{IA}}\,G_{\mathrm{PGA}}
\left[k_d-k_bH_{\mathrm{BP}}(s)\right]H_{\mathrm{LP}}(s),
\label{eq:afe-bloques}
\end{equation}

donde $G_{\mathrm{IA}}$ representa la etapa diferencial, $G_{\mathrm{PGA}}$ la ganancia
seleccionable, $H_{\mathrm{BP}}$ el MFB de compensaci\'on, $k_d$ y $k_b$ las relaciones
del sumador y $H_{\mathrm{LP}}$ el filtro anti-\emph{aliasing}. El orden funcional es:

\begin{center}
\small
\fbox{$v_d$ diferencial}
$\rightarrow$
\fbox{IA/PGA}
$\rightarrow$
\fbox{BP MFB}
$\rightarrow$
\fbox{SUM}
$\rightarrow$
\fbox{LP anal\'ogico}
$\rightarrow$
\fbox{AMux}
$\rightarrow$
\fbox{ADC $\Delta\Sigma$}
\end{center}

La Fig.~\ref{fig:topdesign-analogico} corresponde a la versi\'on consolidada previa al
cambio VDAC$\to$IDAC de agosto. Se conserva porque muestra con claridad la topolog\'ia;
la revisi\'on actual mantiene las cuatro referencias, pero las genera con IDAC y
resistencias externas.

\begin{figure}[ht]
\centering
\includegraphics[width=\textwidth,trim=115 115 0 105,clip]{Schematic-images-0.jpg}
\caption{Cadena de acondicionamiento consolidada en PSoC Creator. El recorte comienza
despu\'es de los terminales de entrada para mantener esta secci\'on independiente del
transductor: PGA diferencial, ganancia adicional, MFB, sumador, LP, referencias
independientes, AMux y ADC. La captura corresponde a la revisi\'on hist\'orica con VDAC.}
\label{fig:topdesign-analogico}
\end{figure}

\begin{figure}[p]
\centering
\includegraphics[width=\textwidth,height=0.78\textheight,keepaspectratio]{AnalogPathGordito.png}
\caption{Lectura funcional ampliada del camino anal\'ogico. Se distinguen el AFE,
las cuatro referencias programables, el multiplexor de observaci\'on y el ADC.
La separaci\'on gr\'afica ayuda a explicar por qu\'e el circuito puede calibrarse
etapa por etapa sin que el lazo permanezca activo durante la adquisici\'on.}
\label{fig:analog-path-gordito}
\end{figure}

\begin{figure}[p]
\centering
\includegraphics[width=\textwidth,height=0.80\textheight,keepaspectratio]{Schematic-Circuit-PSoC.jpg}
\caption{TopDesign completo utilizado como base de la documentaci\'on de la
revisi\'on con VDAC. Adem\'as del AFE aparecen ADC, bloques digitales, DMA,
temporizadores y comunicaciones. Esta captura se mantiene completa para
conservar las relaciones que se pierden en el recorte de la
Fig.~\ref{fig:topdesign-analogico}.}
\label{fig:topdesign-completo}
\end{figure}

Los valores transferidos desde el \texttt{TopDesign.cysch} actual al esquema KiCad se
resumen en la Tabla~\ref{tab:componentes-afe}. La tabla documenta la realizaci\'on f\'isica;
no reemplaza la verificaci\'on de la funci\'on de transferencia completa con tolerancias.

\begin{table}[ht]
\centering
\footnotesize
\caption{Redes externas de la cadena anal\'ogica vigente.}
\label{tab:componentes-afe}
\begin{tabularx}{\textwidth}{@{}p{0.22\textwidth} Y@{}}
\toprule
\textbf{Bloque} & \textbf{Componentes identificados} \\
\midrule
Polarizaci\'on de entrada & Dos resistencias de \SI{50}{\kilo\ohm} hacia la referencia com\'un. \\
Pasa-banda MFB, rama de entrada & $R_4=\SI{43}{\kilo\ohm}$ en serie con $C_1=\SI{680}{\micro\farad}$ (electrol\'itico; acopla en AC y bloquea continua, no fija la banda). \\
Pasa-banda MFB, rama de realimentaci\'on & $R_5=\SI{47}{\kilo\ohm}$ en paralelo con $C_2=\SI{27}{\pico\farad}$ (define el polo de alta frecuencia de esta etapa). \\
Sumador & Entrada directa \SI{6.8}{\kilo\ohm}; potenci\'ometro de \SI{2}{\kilo\ohm} al 31,7\,\% para la rama BP; realimentaci\'on \SI{27}{\kilo\ohm}$\parallel$\SI{15}{\nano\farad}. \\
Pasa-bajos & Entrada \SI{12}{\kilo\ohm}; realimentaci\'on \SI{150}{\kilo\ohm}$\parallel$\SI{3.3}{\nano\farad}; \SI{47}{\nano\farad} de desacoplo asociado. \\
Referencias de etapa & Cuatro ramas de \SI{30}{\kilo\ohm}, desacopladas con \SI{1}{\micro\farad}$\parallel$\SI{100}{\nano\farad}; excitaci\'on IDAC en la rama vigente. \\
AMux & Capacitor externo de \SI{100}{\nano\farad} y selecci\'on de PGA, BP, SUM o LP para observaci\'on. \\
\bottomrule
\end{tabularx}
\end{table}

\begin{figure}[p]
\centering
\includegraphics[width=\textwidth,height=0.82\textheight,keepaspectratio]{kicad_esquema_completo.png}
\caption{Esquema el\'ectrico completo trasladado a KiCad: carrier del PSoC y el
ESP32, perif\'ericos, entrada del ge\'ofono, redes externas del pasa-banda y
sumador, pasa-bajos, referencias y desacoplos. La imagen es deliberadamente
extensa; sirve como plano de implementaci\'on y como reserva para recortes
posteriores.}
\label{fig:kicad-esquema-completo}
\end{figure}

\subsection{Verificaci\'on cuantitativa a partir de los valores reales}
\label{sec:verificacion-cuantitativa-afe}

La Tabla~\ref{tab:componentes-afe} no es una lista decorativa: sus valores permiten
comprobar si el circuito realmente construido hace lo que enuncia la
Ec.~\eqref{eq:comp-afe}, en lugar de asumirlo. Con las resistencias y capacitores del
esquema KiCad vigente,

\begin{equation}
f_{\mathrm{in}}=\frac{1}{2\pi R_4 C_1}\approx\SI{5.4}{\milli\hertz},
\qquad
f_{\mathrm{fb}}=\frac{1}{2\pi R_5 C_2}\approx\SI{125}{\kilo\hertz}.
\label{eq:polos-mfb-reales}
\end{equation}

Ambos polos caen muy lejos de la banda de inter\'es (\SI{10}{}--\SI{50}{\hertz}, la
banda \'util de campo identificada en el historial de mediciones). Esto significa que,
en esa banda, $C_1$ se comporta como un cortocircuito ($X_{C_1}\ll R_4$) y $C_2$ como un
circuito abierto ($X_{C_2}\gg R_5$): la etapa opera con ganancia esencialmente resistiva

\begin{equation}
\left.\frac{v_{\mathrm{BP}}}{v_{\mathrm{in}}}\right|_{\SI{10}{}-\SI{50}{\hertz}}
\approx-\frac{R_5}{R_4}=-\frac{47}{43}\approx-1{,}09.
\label{eq:ganancia-bp-real}
\end{equation}

En otras palabras, \textbf{la rama MFB, tal como est\'a poblada, no realiza un
pasa-banda resonante dentro del rango \'util}: act\'ua como un inversor de ganancia
casi unitaria, acoplado en AC por $C_1$ y limitado en ancho de banda muy por encima de
la zona de trabajo por $C_2$. La forma resonante de la Ec.~\eqref{eq:comp-afe},
con $\omega_n$ y los factores de amortiguamiento $\zeta_0,\zeta_1$, describe el
\emph{objetivo} de compensaci\'on frente al polo mec\'anico del ge\'ofono; los valores
poblados en esta revisi\'on de hardware no ubican ese par de polos/cero dentro de la
banda de campo. Dos lecturas son posibles y ninguna puede descartarse sin medir la
placa real: (a) $R_4/C_1$ y $R_5/C_2$ cumplen aqu\'i s\'olo funciones de acoplo AC y
limitaci\'on de ancho de banda, y la compensaci\'on resonante depende de otro par de
componentes que este inventario no aisl\'o correctamente a partir del esquem\'atico, o
(b) la topolog\'ia poblada en agosto qued\'o, de hecho, m\'as cerca de un inversor
ancho de banda que de un pasa-banda estrecho, y la compensaci\'on de
Ec.~\eqref{eq:comp-afe} se apoya principalmente en el sumador y el pasa-bajos. Esta
verificaci\'on se agrega a la lista de pendientes de la
Secci\'on~\ref{sec:afe-consolidada} y debe resolverse con una medici\'on de magnitud y
fase extremo a extremo, no s\'olo con el an\'alisis de componentes.

El sumador y el pasa-bajos, en cambio, s\'i quedan caracterizados de forma directa
porque cada uno usa un solo capacitor de realimentaci\'on (topolog\'ia inversora de
un polo):

\begin{equation}
f_{\mathrm{SUM}}=\frac{1}{2\pi R_8 C_3}\approx\SI{393}{\hertz},
\qquad
f_{\mathrm{LP}}=\frac{1}{2\pi R_{10} C_5}\approx\SI{322}{\hertz}.
\label{eq:polos-sum-lp}
\end{equation}

Ambos polos quedan muy por encima de la banda \'util de \SI{10}{}--\SI{50}{\hertz} y
por debajo de la frecuencia de Nyquist nativa ($\SI{2604}{SPS}/2\approx\SI{1302}{\hertz}$),
de modo que aportan atenuaci\'on anti-\emph{aliasing} adicional sin recortar la se\~nal
s\'ismica. Las ganancias resistivas de continua tambi\'en se desprenden directamente de
la tabla: el camino directo del sumador aporta $-R_8/R_6\approx-3{,}97$, la rama BP
contribuye a trav\'es del potenci\'ometro con $-R_8/R_{\mathrm{pot}}\approx-42{,}6$ (con
el cursor al 31,7\,\%, es decir $R_{\mathrm{pot}}\approx\SI{634}{\ohm}$) y el pasa-bajos
aporta $-R_{10}/R_9=-12{,}5$. Estos n\'umeros no reemplazan
la medici\'on de magnitud y fase de las Figs.~\ref{fig:analog-path-check},
\ref{fig:medicion-pga-adc} y~\ref{fig:medicion-geo-adc}, pero permiten distinguir, antes
de medir, qu\'e desviaciones observadas se explican por los componentes poblados y
cu\'ales exigir\'ian revisar la topolog\'ia.

Finalmente, la descripci\'on del propio esquem\'atico marca un riesgo de \emph{bring-up}
a\'un no cerrado: $C_1$ es un capacitor electrol\'itico polarizado y su documentaci\'on
indica expl\'icitamente que la polaridad debe verificarse al energizar la placa por
primera vez. Una polaridad inversa en ese punto no impedir\'ia necesariamente el
funcionamiento a bajo nivel de se\~nal, pero degradar\'ia la fuga y el ruido del
acoplo AC de forma dif\'icil de distinguir, en una medici\'on r\'apida, de ruido propio
del ADC.

\subsection{Asistencia digital para recentrar la cadena}

El multiplexor, el ADC y las referencias programables convierten al AFE en una
instrumentaci\'on anal\'ogica asistida digitalmente. Durante una calibraci\'on
\emph{foreground}, cada etapa se observa de manera secuencial y se ajusta su referencia.
Durante la captura, los c\'odigos quedan fijos: no hay un lazo digital modulando la se\~nal
s\'ismica.

Para continua, la salida de la etapa $i$ puede modelarse como

\begin{equation}
y_i^{\mathrm{DC}}=\sum_{j<i}A_{ij}y_j^{\mathrm{DC}}+b_i+K_i u_i,
\label{eq:offset-etapas}
\end{equation}

donde $b_i$ es su \emph{offset}, $A_{ij}$ la propagaci\'on de residuos anteriores,
$u_i$ el c\'odigo o tensi\'on de referencia y $K_i$ su efecto sobre la salida. La forma
triangular de la Ec.~\eqref{eq:offset-etapas} impone el orden PGA$\to$BP$\to$SUM$\to$LP:
ajustar primero una etapa posterior dejar\'ia que el residuo de una anterior volviera a
desplazarla.

La evoluci\'on fue manual al comienzo, con c\'odigos de referencia enviados desde la
interfaz de banco, y despu\'es autom\'atica mediante un controlador PI con banda muerta,
saturaci\'on y \emph{anti-windup}. Los valores aceptados se conservan en EEPROM con CRC
e indexaci\'on por configuraci\'on de ganancia. La revisi\'on de agosto reemplaz\'o los VDAC
por IDAC, pero conserva el mismo principio: medir una salida intermedia, corregir su
referencia y retirar el lazo antes de adquirir.

\begin{figure}[p]
\centering
\includegraphics[width=\textwidth,height=0.80\textheight,keepaspectratio]{psoc_foreground_calibration_flow.drawio_horizontal_comprimido.png}
\caption{Flujo de calibraci\'on \emph{foreground}. Al iniciar o cambiar ganancia
se intenta recuperar un registro EEPROM; si no es v\'alido, el AMux observa PGA,
BP, SUM y LP de manera secuencial, se ajustan las referencias y luego se restaura
el camino de adquisici\'on con los c\'odigos congelados.}
\label{fig:flujo-calibracion-foreground}
\end{figure}

\begin{figure}[p]
\centering
\includegraphics[width=\textwidth,height=0.80\textheight,keepaspectratio]{autocalibration_stage_outputs_5_to_20s.png}
\caption{Salidas internas observadas durante una secuencia de autocalibraci\'on:
PGA, pasa-banda, sumador y pasa-bajos. La figura documenta que el sistema puede
medir nodos intermedios mediante el mismo ADC; no demuestra por s\'i sola que la
calibraci\'on sea repetible entre placas.}
\label{fig:salidas-autocalibracion}
\end{figure}

\begin{figure}[p]
\centering
\includegraphics[width=\textwidth,height=0.80\textheight,keepaspectratio]{autocalibration_filter_comparison.png}
\caption{Comparaci\'on de una configuraci\'on FIR inicial y otra de mayor longitud
durante el desarrollo de la autocalibraci\'on. La imagen permite discutir el
compromiso entre suavizado, respuesta transitoria, costo del DFB y preservaci\'on
de la forma de onda.}
\label{fig:comparacion-fir-autocalibracion}
\end{figure}

\subsection{Digitalizaci\'on: ADC, rango y cuantizaci\'on}

El conversor seleccionado dentro del PSoC~5LP es un $\Delta\Sigma$ configurado a
\textbf{18 bits nominales}. Esta arquitectura resulta apropiada para priorizar resoluci\'on
en una banda relativamente baja; los 18 bits describen la palabra nominal, no garantizan
18 bits efectivos en presencia de ruido, no linealidad y referencias reales
~\cite{InfineonPSoC}.

El firmware vigente emplea una tasa nativa de

\begin{equation}
F_{s,0}=\SI{2604}{SPS}.
\end{equation}

El ADC posee cuatro configuraciones de rango seleccionables. Aplicando
\eqref{eq:adc-cuantizacion}, sus pasos ideales son:

\begin{table}[ht]
\centering
\footnotesize
\caption{Rangos vigentes del ADC $\Delta\Sigma$ de 18 bits.}
\label{tab:rangos-adc}
\begin{tabular}{@{}clrr@{}}
\toprule
\textbf{C\'odigo} & \textbf{Rango bipolar} & $V_{\mathrm{span}}$ & $\Delta V_{\mathrm{LSB}}$ ideal \\
\midrule
1 & $\pm\SI{2.5}{\volt}$   & \SI{5.000}{\volt} & \SI{19.1}{\micro\volt} \\
2 & $\pm\SI{0.512}{\volt}$ & \SI{1.024}{\volt} & \SI{3.91}{\micro\volt} \\
3 & $\pm\SI{1.024}{\volt}$ & \SI{2.048}{\volt} & \SI{7.81}{\micro\volt} \\
4 & $\pm\SI{0.625}{\volt}$ & \SI{1.250}{\volt} & \SI{4.77}{\micro\volt} \\
\bottomrule
\end{tabular}
\end{table}

Elegir el rango es una segunda forma de adaptar la se\~nal al conversor. Pasar de
$\pm\SI{2.5}{\volt}$ a $\pm\SI{0.512}{\volt}$ reduce el LSB ideal casi cinco veces,
pero reduce en la misma proporci\'on el margen antes del recorte. Por tanto, PGA y rango
del ADC deben configurarse conjuntamente: \textbf{la mejor resoluci\'on no es el rango m\'as
estrecho, sino el rango m\'as estrecho que no sature la captura}.

\subsection{Del ADC a memoria: FIR, DMA y control determinista}

La primera digitalizaci\'on utilizaba una interrupci\'on convencional por muestra y salida
serie. El proyecto consolidado evolucion\'o hacia dos caminos concurrentes:

\begin{align}
\text{crudo:}\quad &\mathrm{ADC}\rightarrow\mathrm{DMA}\rightarrow\mathrm{RAM},\\
\text{filtrado:}\quad &\mathrm{ADC}\rightarrow\mathrm{DFB/FIR}
\rightarrow\mathrm{DMA}\rightarrow\mathrm{RAM}.
\end{align}

El bloque DFB ejecuta el FIR en hardware y los canales DMA trasladan las palabras sin que
el Cortex--M3 copie cada muestra. Esto reduce carga de CPU y, sobre todo, separa el
instante de muestreo de tareas lentas de control o diagn\'ostico.

\begin{figure}[ht]
\centering
\includegraphics[width=\textwidth]{DigitalPathGordito_cropped.png}
\caption{Camino digital del proyecto consolidado: filtro DFB, l\'ogica de control,
solicitudes DMA, registros de estado, EEPROM y temporizadores. La UART que aparece en
esta captura corresponde a la revisi\'on anterior; en la rama de hardware vigente la
salida del PSoC se migr\'o a I$^2$C y la UART qued\'o s\'olo como recepci\'on.}
\label{fig:camino-digital-nueva}
\end{figure}

\begin{figure}[p]
\centering
\includegraphics[width=\textwidth,height=0.82\textheight,keepaspectratio]{system_architecture_autocalibration.drawio.png}
\caption{Arquitectura conjunta de adquisici\'on y autocalibraci\'on. La frontera
roja delimita el hardware anal\'ogico, la verde los recursos digitales del PSoC
y la violeta la comunicaci\'on mediante ESP32. El diagrama muestra d\'onde termina
esta secci\'on y d\'onde comienza la transmisi\'on.}
\label{fig:arquitectura-autocalibracion}
\end{figure}

\begin{figure}[p]
\centering
\includegraphics[width=\textwidth,height=0.80\textheight,keepaspectratio]{blocks_topdesign_schematic.png}
\caption{TopDesign reorganizado por dominios funcionales: frente anal\'ogico,
referencias, ADC/AMux, entradas auxiliares, gesti\'on de eventos, maestro ESP32 y
esclavo. Las flechas resaltan las fronteras de datos, control y sincronizaci\'on.}
\label{fig:topdesign-por-bloques}
\end{figure}

Desde julio, el bloque \texttt{superMaquina}, descrito en Verilog e implementado en UDB,
decide en hardware cu\'ando habilitar cada solicitud DMA, cu\'ando comienza y termina una
captura y qu\'e eventos pueden interrumpir al procesador. La regla de partici\'on es:

\begin{itemize}
    \item va a \textbf{hardware} lo que debe ocurrir en un instante preciso: flanco de
    inicio, habilitaci\'on del ADC, conteo de muestras/lotes, enrutamiento DMA y fin;
    \item queda en \textbf{software} lo lento pero complejo: calibraci\'on, persistencia,
    configuraci\'on, supervisi\'on y protocolo de comandos.
\end{itemize}

El sistema admite decimaci\'on programable

\begin{equation}
F_s=\frac{\SI{2604}{SPS}}{N},\qquad N\in\{1,\ldots,100\}.
\label{eq:decimacion-psoc}
\end{equation}

El promediado de $N$ muestras reduce ruido no correlacionado, pero su respuesta es una
\emph{sinc}; no constituye por s\'i solo un filtro anti-\emph{aliasing} ideal. Por eso la
capacidad de seleccionar $N$ no implica que los cien valores sean igualmente apropiados:
para cada configuraci\'on de campo debe verificarse la banda combinada del LP anal\'ogico,
el FIR y la decimaci\'on.

\begin{figure}[p]
\centering
\includegraphics[width=\textwidth,height=0.80\textheight,keepaspectratio]{medicion_pga_adc.png}
\caption{Comparaci\'on exploratoria PGA$\rightarrow$ADC: magnitud, fase y
coherencia frente al modelo y a envolventes de tolerancia. La coherencia alta en
la banda baja respalda el procedimiento de identificaci\'on; su deterioro fuera
de esa banda marca el l\'imite de interpretaci\'on de la medici\'on.}
\label{fig:medicion-pga-adc}
\end{figure}

\begin{figure}[p]
\centering
\includegraphics[width=\textwidth,height=0.80\textheight,keepaspectratio]{medicion_geo_adc.png}
\caption{Comparaci\'on exploratoria GEO$\rightarrow$ADC, incorporando el modelo
del ge\'ofono a la cadena electr\'onica. La separaci\'on entre puntos crudos y
procesados a frecuencias altas refuerza la necesidad de no extrapolar una
identificaci\'on fuera de la regi\'on con coherencia suficiente.}
\label{fig:medicion-geo-adc}
\end{figure}

\clearpage

\subsection{Hitos dentro del proyecto consolidado}

Una vez absorbidos los proyectos exploratorios, las modificaciones dejaron de crear
carpetas nuevas y pasaron a actualizar \texttt{AcondicionamientoAnalogico}. Los cambios
estructurales principales recuperados con Git son:

\begin{table}[ht]
\centering
\footnotesize
\begin{tabularx}{\textwidth}{@{}p{0.13\textwidth} p{0.13\textwidth} Y@{}}
\toprule
\textbf{Fecha} & \textbf{Commit} & \textbf{Cambio de arquitectura} \\
\midrule
03/06 & \texttt{4062297} & Nace el proyecto unificado con la cadena completa. \\
17/06 & \texttt{e490334} & Se consolidan DMA, camino filtrado por DFB, EEPROM y selecci\'on crudo/FIR; se eliminan del \emph{head} los proyectos exploratorios, que permanecen recuperables en Git. \\
02/07 & \texttt{87dcd7e} & La l\'ogica temporal de captura migra a \texttt{superMaquina} en Verilog/UDB. \\
11/07 & \texttt{38098e9} & Se consolida la tasa nativa de \SI{2604}{SPS} y la decimaci\'on asociada. \\
03/08 & \texttt{f355a8a} & Nueva revisi\'on de hardware: cuatro IDAC sustituyen los VDAC, se agrega una etapa \texttt{PGAout}, la salida de datos pasa a I$^2$C y la UART queda s\'olo RX. \\
\bottomrule
\end{tabularx}
\end{table}

\subsection{Qu\'e queda establecido y qu\'e debe verificarse}

La arquitectura ya est\'a definida: entrada diferencial, ganancia programable,
compensaci\'on MFB con sumador, pasa-bajos anal\'ogico, referencias independientes,
observaci\'on por AMux, ADC $\Delta\Sigma$ de 18 bits, cuatro rangos, DFB/FIR, DMA y
control temporal en hardware. Git permite adem\'as justificar esa arquitectura como el
resultado de alternativas concretas, no como una descripci\'on retrospectiva.

Antes de cerrar esta secci\'on como dise\~no validado faltan pruebas que pertenecen
exclusivamente a la instrumentaci\'on:

\begin{itemize}
    \item respuesta en magnitud y fase extremo a extremo para cada ganancia y rango;
    \item ruido referido a la entrada y n\'umero efectivo de bits, no s\'olo resoluci\'on
    nominal;
    \item tolerancias entre varias placas y deriva t\'ermica;
    \item margen de saturaci\'on en cada nodo intermedio, no \'unicamente en la salida;
    \item rechazo anti-\emph{aliasing} para las combinaciones de FIR y decimaci\'on que se
    usar\'an;
    \item repetibilidad de la calibraci\'on luego de la migraci\'on VDAC$\to$IDAC;
    \item confirmaci\'on por medici\'on de qu\'e etapa realiza efectivamente la
    compensaci\'on resonante de la Ec.~\eqref{eq:comp-afe}, dado que los valores
    poblados de la rama MFB (Secci\'on~\ref{sec:verificacion-cuantitativa-afe}) no
    ubican sus polos dentro de la banda de campo;
    \item polaridad de $C_1$ (electrol\'itico de acoplo AC en la rama de entrada del
    MFB), marcada como pendiente de revisi\'on desde el propio esquem\'atico.
\end{itemize}

Estas verificaciones no cambian el alcance de la etapa: determinan si el camino entre
los terminales de entrada y la memoria digital cumple lo dise\~nado. La transmisi\'on de
los datos almacenados comienza reci\'en despu\'es de este punto.
